\section{Применение многосторонних конфиденциальных вычислений в современных системах}
\label{sec:applications}

\subsection{Приватная аналитика данных и аукционы}

В эстонском проекте Sharemind Министерство образования и науки и Налогово-таможенный департамент распределили сведения об обучении и налоговых выплатах между тремя серверами и совместно исследовали связь занятости студентов с завершением обучения. Согласно описанию Cybernetica от 7 июля 2021 г., анализ выполнялся без раскрытия индивидуальных записей и не выявил связи между работой во время учёбы и несвоевременным окончанием вуза, но показал зависимость между уровнем образования и доходом \cite{18}.

\subsection{Финансовые приложения и управление ключами}

При распределённом управлении ключом криптовалютного кошелька закрытый ключ разделяется между несколькими участниками, а подпись транзакции формируется совместно без восстановления секрета в одном месте. Такая пороговая схема устраняет единую точку компрометации и позволяет заранее определить число участников, необходимое для подтверждения операции \cite{06}.

\subsection{Машинное обучение и анализ больших данных}
\label{subsec:ml-applications}

Общая характеристика. МО часто требует объединения наборов данных; МКВ позволяет обучать модели и выполнять инференс без передачи данных в незащищённом виде. Evans et al. описывают облачное инференс-обслуживание и совместное обучение. В первом случае протокол объединяет скрытый запрос пользователя и параметры модели сервера, возвращая предсказание только пользователю. Во втором несколько организаций обучают модель на объединённых данных без доступа к исходным наборам друг друга, используя секретные доли и часто гибридные методы МКВ и гомоморфного шифрования.

Платформы SecureML, Crypten и Cerebro. Обзор \cite{07} связывает интеграцию МКВ в МО с доступностью библиотек. CRYPTEN исследовательского подразделения Facebook автоматически преобразует модели PyTorch в распределённую МКВ-\allowbreak версию, поддерживает GPU-\allowbreak ускорение и semi-\allowbreak honest модель угроз. Интерфейс, имитирующий API стандартных ML-библиотек, упрощает внедрение. Cerebro реализует доменно-специфический язык описания алгоритмов обучения, оптимизатор и механизмы аудита; компилятор выбирает наиболее эффективные МКВ-протоколы. Механизмы \textit{compute policies} и криптографического аудита контролируют выпуск результатов и корректность вычислений без раскрытия данных.

Обучение с двумя серверами (SecureML).  В классической работе Mohassel \& Zhang (SecureML, 2017) описана схема двухсерверного обучения, при которой данные распределяются между двумя некоррумпированными серверами.  Протокол строит модель линейной или логистической регрессии с использованием стохастического градиентного спуска: в каждой итерации веса \(w_j\) обновляются по правилу

\begin{equation}
w_j := w_j - \alpha \cdot \frac{\partial C_i(w)}{\partial w_j},
\end{equation}

где \(\alpha\) --- шаг обучения, \(C_i\) --- функция потерь для выбранного образца, а \(\partial C_i(w) / \partial w_j\) --- частная производная по параметру \(w_j\).  Для логистической регрессии поверх линейного слоя добавляется сигмоидальная функция активации; нейронные сети строятся как обобщение этих слоёв.  Все арифметические операции выполняются на секретных долях, а сложные функции (например, сигмоидальная активация) аппроксимируются полиномами для сокращения затрат на коммуникацию.  SecureML исключает необходимость доверенного посредника, поскольку данные остаются разделёнными между серверами.

Протоколы глубокого обучения на основе Shamir-\allowbreak деления.  Исследование Zhang et al. (2025) \cite{12} показывает, что для глубокой нейронной сети эффективнее использовать пороговое Shamir-деление.  Авторы предлагают схему с фиксированной точкой для представления вещественных чисел и алгоритм усечения, который ограничивает рост десятичных знаков при умножении; это обеспечивает корректные операции свёртки и функции Softmax, достигая точности более 99 \% при нескольких итерациях.  Предлагаемый протокол оптимизирует умножения с помощью алгоритма Винограда, что снижает число мультипликативных врат более чем на 50 \%.  Основным преимуществом Shamir-базированных схем является более высокая эффективность по сравнению с аддитивным делением, особенно на этапе подготовки и онлайновом умножении; однако поддержание точности требует специальных процедур усечения и поддержки фиксированной запятой.

Проблемы и направления. Ограничения коммуникации, аппроксимации нелинейных функций и глубины схем из-за роста шума рассмотрены в разделе~\ref{sec:limitations}. Разработка библиотек направлена на автоматический подбор протоколов и GPU-ускорение (см. описание платформ выше).

\subsection{Медицина и биоинформатика}

При совместном обучении диагностической модели клиниками данные пациентов остаются локальными, а МКВ защищает вычисление агрегированных обновлений. Обзор \cite{11} относит этот сценарий к основным направлениям конфиденциального медицинского МО; применимость определяется приватностью, вычислительными издержками и масштабируемостью.

\subsection{Блокчейн и децентрализованные приложения}

В протоколе DHSМКВ, предложенном Bao et al. для блокчейн- и облачных систем, многоключевое гомоморфное шифрование позволяет выполнять вычисления над шифротекстами, а направленная расшифровка открывает результат только авторизованной стороне. Подход защищает данные и промежуточные значения без требования взаимного доверия участников \cite{05}.

\subsection{Прочие применения: голосование, доказательства с нулевым разглашением и мониторинг}

В электронном голосовании МКВ позволяет конфиденциально суммировать бюллетени; в церемониях формирования параметров доказательств с нулевым разглашением --- распределять генерацию секрета так, чтобы ни один участник не владел им полностью; в сетевом мониторинге --- совместно выявлять совпадения и аномалии без объединения открытых журналов. Во всех сценариях раскрывается только согласованный результат, а исходные данные остаются закрытыми \cite{06}.

\subsection{Вывод}

В следующем разделе рассмотрены ограничения внедрения МКВ и способы их преодоления.